\subsection{A general inversion obstruction from the $C_3+3$ TPE}
\begin{proposition}
Every complex near-group category of type $(C_3,3)$ is Galois conjugate to
a unitary category. Consequently $R8\_3\_436$ has no complex categorification.
\end{proposition}
\begin{proof}
Thornton's sphericality and Galois-pseudo-unitarity theorem
\cite[Theorem~1.3]{ThorntonNear} allows us to take the noninvertible
dimension $d=(3+\sqrt{21})/2$. Its second indicator is $+1$ by
\cite[Theorem~3.9]{LPRV}, since its self-multiplicity is odd.
We may therefore use the complete spherical tetrahedral parametrization
of \cite[Theorem~2.3, Proposition~2.7, Theorem~4.12]{NearGroupTPE2026}.

Write $w^2+w+1=0$, $q=1+w$, and $t=-1/d=(3-d)/3$.
Both possible nondegenerate bicharacters are $\chi(g,h)=w^{gh}$;
the even refinement is forced to be $a(0)=1$, $a(1)=a(2)=w$.
Put $x=\beta(1)$ and $y=\beta(2)$.
The product identity and two cubic equations give
\[
xy=w/3,\qquad y^2+qtx=tw,\qquad x^2+qty=tw.
\]
Thus $(y-x)(x+y-qt)=0$. The case $x=y$ would imply
$6t^2-6t+1=0$, contrary to $t^2-t=1/3$.
Hence $x+y=qt$, $xy=w/3$, and the rotation scalar is uniquely
$c=1/(2+w)$. The exact verifier checks that these relations are equivalent
to all nine cubic equations, including both dimension roots and both
bicharacters.

The numbers $u=x/q$ and $v=y/q$ are roots of the real polynomial
$z^2-tz+1/3$, whose discriminant is $-d/3<0$.
Therefore $v=\overline u$, so
$\overline x=y/w$, $\overline y=x/w$.
These are the Hermitian scalar conditions, and
\cite[Corollary~5.5]{NearGroupTPE2026} unitarizes the complete table of
the original category. The completeness of that table is essential;
an existence-only reconstruction would not justify this conclusion.

For the source tensor, the indices $[0,1,2,5]$ give this near-group ring,
and the invertible at index $3$ conjugates indices $1$ and $2$.
Galois conjugation transports this ordinary tensor autoequivalence along
with the category, preserving its label permutation. The unitary
classification \cite[Theorem~9.1, Example~9.4]{IzumiNear} identifies
the resulting neutral category with an Izumi--Xu category.
Both have trivial outer group of ordinary complex tensor
autoequivalences \cite[Section~4.2, p.~10]{GJSCyclic}; inner
autoequivalences fix labels since the ring is commutative.
The required inversion is impossible.
\end{proof}
This new argument uses the full spherical parametrization in the October
2026 preprint. The exact checker verifies all source coefficients and the
polynomial reduction; the cited categorical inputs and the new argument
remain subject to mathematical review.
